\section{Live experimentation and the measured breakthrough}
\label{sec:live-method}

\subsection{Three scales that must not be conflated}
The account's latest Elo is a stateful matchmaking variable. A deployment's
performance rating estimates a constant strength against the opponents
actually faced. A randomized arm contrast estimates the effect of changing a
component while the account, time period and surrounding software are shared.
They answer different questions. Figure~\ref{fig:modern-performance} displays
performance segments; Table~\ref{tab:c13-arms} gives the controlled contrast.

For outcomes $s_i\in\{0,\tfrac12,1\}$ and known pre-game opponent ratings
$O_i$, \code{search/elo_monitor.py} solves
\begin{equation}
  \sum_i \frac{1}{1+10^{(O_i-\widehat R)/400}}=\sum_i s_i.
  \label{eq:performance-rating}
\end{equation}
Under independent Bernoulli outcomes and a fixed Elo logistic model, let
$c=\log(10)/400$ and $p_i=1/(1+10^{(O_i-\widehat R)/400})$. The log-likelihood
score is $c\sum_i(s_i-p_i)$ and the information is
$c^2\sum_ip_i(1-p_i)$. The implementation's approximate interval half-width
is therefore
\begin{equation}
  1.96\,\frac{400/\log(10)}{\sqrt{\sum_i p_i(1-p_i)}}.
  \label{eq:performance-interval}
\end{equation}
Tie scores use the same estimating equation rather than a separately fitted
draw model. All-wins/all-losses samples clamp to the supported range and have
no finite ordinary interval. Opponent skill drift, repeated opponents and
serial dependence make these intervals descriptive. They are neither
sequential confidence sequences nor independent replications of a prior gate.

\begin{figure}[htbp]
\centering
\begin{tikzpicture}
\begin{axis}[
 width=.70\linewidth,height=7.0cm,
 y dir=reverse,ytick={1,2,3,4,5},ymin=.4,ymax=5.6,
 yticklabels={{Early search (1325)},{C4 preview (635)},
 {C5 plus browser repair (1169)},{Reconstruction v2 / rain (992)},
 {Mev after C13 (455)}},
 yticklabel style={font=\scriptsize,text=ink},
 xmin=1170,xmax=1400,xtick={1200,1250,1300,1350,1400},
 xlabel={\small performance Elo with descriptive approximate 95\% interval},
 ticklabel style={font=\scriptsize,text=muted},
 xmajorgrids,grid style={draw=gridgray},axis line style={draw=muted},clip=false]
 \addplot[only marks,mark=*,mark size=2.6pt,color=seriesA,
 error bars/.cd,x dir=both,x explicit] coordinates {
 (1205.94,1) +- (19.44,0)
 (1236.27,2) +- (28.23,0)
 (1287.48,3) +- (20.63,0)
 (1235.55,4) +- (22.51,0)};
 \addplot[only marks,mark=diamond*,mark size=3.3pt,color=seriesB,
 error bars/.cd,x dir=both,x explicit] coordinates {
 (1343.72,5) +- (33.33,0)};
\end{axis}
\end{tikzpicture}
\caption{Chronological performance cohorts, with sample size in parentheses.
The Mev cohort joins all completed C13final assignments across the permanent
team-setting restart. The earlier rows are contiguous non-A/B policy segments
6, 7, 8 and 13 in the frozen monitor snapshot. Time, opponents, repairs and
policy differ across rows: this chart is not a causal decomposition.}
\label{fig:modern-performance}
\end{figure}

\subsection{Balanced live assignment and the sequential rule}
C9--C14 use the same interleaved A/B infrastructure. A frozen JSON spec names
base settings and binary factor overrides. With $k$ unfrozen factors, a
seeded permutation of all $2^k$ combinations fills each balanced block.
Assignment is stored by room before play; a resumed room reuses its arm.
A one-factor comparison alternates the two treatments in seeded two-game
blocks. Both arms share one account and matchmaking history. This makes
contemporaneous opponent populations comparable in expectation, without
claiming that each arm faced identical human opponents or that time and
rating feedback disappear entirely.

Let $\bar s_1,\bar s_0$ be on/off scores, $d=\bar s_1-\bar s_0$, and
$S^2$ the pooled sample variance across all on/off game scores. The code uses
$V=S^2(1/n_1+1/n_0)$ and the transform
\begin{equation}
  d_j=\frac{1}{1+10^{-E_j/400}}-\frac12,\qquad
  \Lambda=\frac{d_1-d_0}{V}
    \left(d-\frac{d_0+d_1}{2}\right),
  \quad (E_0,E_1)=(0,40).
  \label{eq:live-llr}
\end{equation}
The nominal Wald boundaries \citep{wald1945} are
\begin{equation}
  \log\frac{\beta}{1-\alpha}=-2.2513,\qquad
  \log\frac{1-\beta}{\alpha}=2.8904,
  \quad\alpha=.05,\ \beta=.10.
\end{equation}
No decision is allowed below 300 completed games. An upper crossing accepts
H1, a lower crossing H0; a 3,000-game cap is inconclusive and keeps the
incumbent. The implemented test is a normal approximation with estimated
variance, minimum sample and cap. It does not inherit exact anytime error
control merely by using Wald's boundaries. Dependence, factorial interactions,
multiple factors and sequential candidate selection further qualify the
nominal error interpretation. Four C9 factors at nominal .05 each were not
a family-wise .05 experiment; neither were C13's two successive team tests.

The displayed Elo effect is
\begin{equation}
  \widehat E=400\log_{10}\frac{.5+d}{.5-d}.
  \label{eq:effect-transform}
\end{equation}
It maps a score difference to an Elo advantage against an equal baseline.
It is not exactly the difference between two independently fitted performance
ratings, although they are close for C13. A fixed-sample interval obtained
by transforming $d\pm1.96\sqrt V$ is descriptive after optional stopping.
No such interval replaces the actual stopping rule.

\paragraph{Outcome and missingness policy.}
Forfeits and opponent timer losses are genuine terminal wins and count.
Excluding them would condition on a consequence of good positions as well as
on opponent behaviour. An assigned room without a terminal record is
excluded until it finishes; it is not scored as a tie. An outdated module
docstring says unfinished games count .5, but the actual \code{join} consumes
only completed terminal games. The frozen article audit follows that code.
A future audit should continue reporting missing assignments by arm, because
terminal-only analysis can still be biased by treatment-dependent outages.
Real ties count .5; none occurred in the frozen 7,734-game live dataset.

\subsection{C13: the team comparison that passed}
The incumbent was \code{Rain-Recife-special-stat-fix}: Garchomp, Politoed,
Incineroar, Archaludon, Farigiraf and Golisopod. Round one tried
\code{Rain-Gengar-Helm-Baltimore}; round two tried
\code{SunRain-Mev-Brisbane} against the surviving rain champion. Both arms
used the same C5 search and C4 pvalue preview. Only the exact six sets changed.

\begin{table}[htbp]
\centering\small
\begin{tabular}{lrrrrr}
\toprule
C13 arm & Games & Wins & Score & Performance & Mean opponent \\
\midrule
R1 Helm & 151 & 66 & .4371 & $1132\pm57$ & 1178.9 \\
R1 Recife & 151 & 82 & .5430 & $1204\pm57$ & 1172.6 \\
\midrule
R2 Mev & 448 & 249 & .5558 & $1286\pm34$ & 1243.5 \\
R2 Recife & 449 & 210 & .4677 & $1225\pm34$ & 1249.6 \\
\bottomrule
\end{tabular}
\caption{First-stopping C13 cohorts, recomputed by joining completed games to
saved room assignments. Later arrivals remain in the audit ledger but do not
change the 302- and 897-game decisions. Intervals are approximate descriptive
95\% performance intervals, not sequential arm-effect intervals.}
\label{tab:c13-arms}
\end{table}

Helm stopped H0 at 302 games: $d=-.10596$, estimated $-74.8$ Elo,
$\Lambda=-2.323$. Mev stopped H1 at 897 games:
\begin{align*}
 d&=249/448-210/449=.0880976,\\
 \sqrt V&\approx.0334,\qquad \widehat E=61.9,\qquad \Lambda=3.054>2.890.
\end{align*}
The independent performance-rating difference is 61.5, close to the
score-transform estimate. Mean opposing ratings differed by only 6.1 points,
with the challenger facing slightly lower opposition. The direct randomized
contrast is stronger evidence than comparing two peaks or entire deployment
periods. It supports the Mev team under this fixed preview and turn pilot;
it does not establish universal team quality or isolate any one species.
The effect estimate was selected at a successful stopping time and may
regress; further independent ladder evidence remains valuable.

After the decision on 30 October, \code{c13final} played Mev every game.
The frozen cohort contains 455 completed games, 226 wins and 229 losses,
performance $1343.7\pm33.3$, against mean opposition 1346.9. Its descriptive
lower interval endpoint is 1310.4. The 49.7\% win rate is compatible with
improvement because matchmaking supplied stronger opponents. This cohort
corroborates sustained play above the old plateau, but is a subsequent
observational cohort rather than another randomized treatment comparison.
C14 games are excluded from it.

The runner's default team changed permanently on 31 October. Before that,
run-level metadata still named the rain default while the per-game spec
selected Mev. Accordingly, the article identifies the champion cohort from
\emph{assignment ledgers}, not by filtering the default team name. Two
monitor segments across that restart are combined. This corrects a subtle
measurement trap without changing any campaign result.

\subsection{How the evidence changed the research process}
The early gate measured sampled actor play; search deployment made that
objective mismatched. The next gate used search opponents; imitation brought
some human behaviours closer, but retained serious action-distribution and
strength differences. C9 made target-domain interventions directly testable.
C10 and C12 subsequently required an offline pass merely to qualify for a
fresh live experiment, never as sufficient evidence for automatic live gain.
C13's team-only contrast finally passed that direct ladder test.

The pooled offline rule also addressed statistical power and wasted compute.
At 20\% discordance, a 20,000-campaign simulation found a +4-point gain
passed the intended old two-stage rule about 9.5\% of the time versus 78.7\%
for the pooled 800-pair rule; null pass rates were .06\% and 4.30\%.
These are assumptions-based power simulations, not observed policy gains.
A better rule cannot fix a poor surrogate opponent, as C12's subsequent
rejection and C13's direct test illustrate.

\source{\code{search/ab_test.py} (assignment, terminal join, score transform,
normal-approximation SPRT); \code{search/elo_monitor.py};
\code{scripts/team_ladder.py}; \code{docs/m5-pipeline-review.md};
\code{evidence-20261031.json}, \code{collect_evidence.py}.}
